\documentclass[14pt]{article}
\usepackage{amsmath, amssymb, amsthm}
\usepackage[utf8]{inputenc}
\usepackage[russian]{babel}

\usepackage{extsizes}

\usepackage[left=25mm, top=20mm, right=10mm, bottom=20mm, nohead, nofoot]{geometry}
\hoffset=-5mm
\voffset=5mm

\usepackage{graphicx}
\graphicspath{{./img}}

\begin{document}
	
	\textbf{Лекция 10}
	
	\vspace{2mm}  % небольшой отступ после заголовка
	
	\includegraphics[scale=0.3]{img_2}
	
	\begin{equation}
		U_{\text{внутр}} + \frac{p}{\rho} = \text{const}
	\end{equation}
	
	\begin{equation}
		U_{\text{внутр}} = U_{\text{т}} + \frac{v^{2}}{2} + gh
	\end{equation}
	
	Энтропия:
	\begin{equation}
		I = U_{\text{т}} + pV
	\end{equation}
	
	\begin{equation}
		i = \frac{I}{m} = \hat{U}_{\text{т}} + \frac{pV}{m}, \quad \text{где } \frac{pV}{m} = \frac{p}{\rho}
	\end{equation}
	
	Следовательно,
	\begin{equation}
		i + \frac{v^{2}}{2} + gh = \text{const}
	\end{equation}
	
	\begin{itemize}
		\item Дж.-Томсон: \( v \approx 0 \), \( \Delta h \approx 0 \) \( \Rightarrow \) \( i = \text{const} \)
		\item \( \Delta U_{\text{т}} \approx 0 \), \( \Delta h \approx 0 \) \( \Rightarrow \) \( \frac{v^{2}}{2} + \frac{p}{\rho} = \text{const} \), \( v \uparrow \) \( \Rightarrow \) \( p \downarrow \)
	\end{itemize}
	
	Второе начало термодинамики:
	
	\textbf{Клаузиуса:} Невозможна самопроизвольная\textsuperscript{*} передача тепла\textsuperscript{**} от менее нагретого тела к более нагретому.
	
	\vspace{0.2cm}
	\noindent
	\textsuperscript{*}без изменения в других телах
	
	\noindent
	\textsuperscript{**}независимо от способа теплопередачи
	
	\textbf{Томсона (Планка):} Невозможен циклический процесс единственным результатом которого является совершение работы за счет уменьшения тепловой энергии какого-либо тела.
	\pagebreak
	
	Пример:
	
	\includegraphics[scale=0.3]{img_1}
	
	\begin{equation}
		\Delta U_{0} = 0 \Rightarrow A_{0} = Q = Q_{\text{н}} - Q_{\text{х}}
	\end{equation}
	
	\begin{equation}
		\eta \stackrel{\text{def}}{=} \frac{A}{Q_{\text{н}}} = \frac{Q_{\text{н}} - Q_{\text{х}}}{Q_{\text{н}}}
	\end{equation}
	
	Если \( Q_{\text{х}} = 0 \), то \( \eta = 100\% \) \(\leftarrow \) "вечный двигатель 2го рода"
	
	Доказательство эквивалентности:
	
	1. Клаузиус \( \Rightarrow \) Т. - П. от противного, пусть существует вечный двигатель 2го рода
	
	\includegraphics[scale=0.3]{img_3}
	
	Итог: \(Q \) от тела с \(T_{1}\) перешло к телу с \(T_{2} > T{1} \) противоречие ч.т.д.
	
	2. Т. - П. \( \Rightarrow \) Клаузиус от противного: циклический процесс, в котором тело получает \(Q_{\text{н}}\) от тела с \(T_{\text{н}}\) и отдает \(Q_{\text{х}}\) телу с \(T_{\text{х}}\)
	
	Приведем в тепловой контакт холодильник и нагреватель так, чтобы холодильник отдал тепло \(Q_{\text{х}}\) 
	
	Тогда:
	\begin{equation}
		\begin{cases}
			Q_{\text{н}} = Q_{\text{н}} - Q_{\text{х}}
			\\
			Q_{\text{х}} = 0
			\\
			A = Q_{\text{н}} - Q_{\text{х}}
		\end{cases}
	\end{equation}
	это вечный двигатель второго рода (Perpetuum mobile II) что невозможно ч.т.д.
	
	\textbf{Об обратимости ТД процессов}
	
	Обратимый процесс - такой процесс, что система может вернуться в исходное состояние А без изменения в других телах.
	
	В узком смысле обратимый процесс - такой процесс, что система может вернуться в исходное состояние А через туже последовательность состояний.
	
	\(\forall\) квазистатический процесс обратим в узком смысле
	
	\includegraphics[scale=0.3]{img_4}
	
	Необратимые процессы (примеры):
	\begin{itemize}
		\item химическая реакция
		\item теплопередача
		\item диффузия
		\item расширение газа в вакууме
	\end{itemize}
	
	\textbf{Цикл Карно и Теорема Карно}
	\begin{equation}
		\eta = \frac{Q_{\text{н}} - Q_{\text{х}}}{Q_{\text{н}}} = \frac{A}{Q_{\text{н}}}
	\end{equation}
	\underline{Теорема Карно:}
	\(\eta = \eta(T_{\text{н}}, T_{\text{х}})\). КПД Цикла Карно зависит только от \(T_{\text{н}}\) и \(T_{\text{х}}\) и не зависит от природы рабочего тела и устройства тепловой машины.
\end{document}